% =====================================================
% VISÃO GERAL DA DISCIPLINA — PROGRAMAÇÃO PARA A INTERNET II
% =====================================================
\section{Visão Geral da Unidade Curricular}

A Unidade Curricular de \textbf{Programação para a Internet II} tem como objetivo aprofundar os conhecimentos adquiridos na disciplina antecedente (Programação para a Internet I), evoluindo o desenvolvimento frontend básico em HTML5, CSS3 e JavaScript Vanilla para a construção de \textbf{Aplicações Web Modernas, Reativas, Modulares e Integradas a Serviços Web RESTful}.

Ao longo do semestre, os estudantes irão desenvolver competências avançadas em arquitetura de código client-side, modularização com ES6 Modules, manipulação avançada de estado no cliente, consumo de APIs assíncronas, roteamento no cliente (Single Page Applications — SPAs) e introdução a ferramentas modernas de build e bibliotecas de interface (Vite e React/Vue).

\subsection*{Estrutura do Semestre}

O semestre está organizado em 5 módulos estruturantes de aprendizagem prática:

\begin{enumerate}
    \item \textbf{Retomada e Diagnóstico (Revisão de PI-I)} \\
    Revisão prática de manipulação do DOM, eventos, assincronismo com \texttt{async/await}, Fetch API e armazenamento com \texttt{localStorage}.

    \item \textbf{Arquitetura Modular e ES6 Modules} \\
    Organização profissional de código frontend, separação de responsabilidades em serviços, componentes e manipuladores de eventos.

    \item \textbf{Gestão Avançada de Estado e Formulários Reativos} \\
    Persistência no navegador, validações avançadas no cliente, tratamento rigoroso de erros e experiência do usuário (UX).

    \item \textbf{Consumo de APIs RESTful e Desenvolvimento de SPAs} \\
    Construção de Single Page Applications com roteamento client-side sem recarregamento de página e integração com rotas de backend.

    \item \textbf{Introdução a Ecossistemas Modernos e Projeto Final} \\
    Uso de bundlers (\texttt{Vite}), introdução à componentização reativa e entrega de um projeto Web completo de final de semestre.
\end{enumerate}

% =====================================================
% PIPELINE DA DISCIPLINA
% =====================================================
\begin{figure}[H]
\centering
\resizebox{\linewidth}{!}{%
\begin{tikzpicture}[
    node distance=1.4cm and 1.0cm,
    every node/.style={font=\small},
    etapa/.style={
        rectangle,
        rounded corners=6pt,
        draw=azulTitulo,
        thick,
        minimum width=3.2cm,
        minimum height=1.55cm,
        text width=3.2cm,
        align=center,
        fill=azulTitulo!5
    },
    seta/.style={
        -{Latex[length=2.2mm]},
        thick,
        draw=azulTitulo
    }
]

\node[etapa] (e1) {\textbf{1. Retomada}\\Revisão DOM, Fetch\\e LocalStorage};

\node[etapa, right=of e1] (e2) {\textbf{2. Módulos ES6}\\Arquitetura limpa\\Separação em camadas};

\node[etapa, right=of e2] (e3) {\textbf{3. Estado \& Forms}\\Validação reativa\\Gestão de estado local};

\node[etapa, right=of e3] (e4) {\textbf{4. SPAs \& REST}\\Consumo de APIs\\Roteamento client-side};

\node[etapa, right=of e4] (e5) {\textbf{5. Projeto Final}\\Vite / Frameworks\\Aplicação Web Completa};

\draw[seta] (e1) -- (e2);
\draw[seta] (e2) -- (e3);
\draw[seta] (e3) -- (e4);
\draw[seta] (e4) -- (e5);

\end{tikzpicture}%
}
\caption{Pipeline de aprendizagem da UC Programação para a Internet II.}
\end{figure}

\subsection*{Competências Desenvolvidas}

Ao final da disciplina, o estudante deverá ser capaz de:

\begin{itemize}
    \item Construir interfaces web altamente dinâmicas e responsivas;
    \item Estruturar aplicações JavaScript utilizando arquitetura modular em ES6 Modules;
    \item Consumir e manipular dados de serviços Web RESTful usando \texttt{fetch} e \texttt{async/await};
    \item Gerenciar estado local e persistência no cliente via Web Storage API;
    \item Implementar padrões de projeto de Single Page Application (SPA);
    \item Utilizar ferramentas modernas de ambiente de desenvolvimento web (\texttt{Node.js}, \texttt{npm}, \texttt{Vite});
    \item Publicar e versionar aplicações web profissionais no GitHub e GitHub Pages.
\end{itemize}

\subsection*{Metodologia de Avaliação}

A avaliação será continuada e baseada na entrega de roteiros práticos semanais, avaliações de desempenho prático em laboratório e no desenvolvimento de um Projeto Web Integrador entregue e apresentado ao final do semestre.
